\documentclass{article}
\usepackage{amsmath,amssymb,amsthm}
\usepackage{algorithm}
\usepackage{algorithmic}
\usepackage{hyperref}
\usepackage{enumitem}
\newtheorem{theorem}{Theorem}
\newtheorem{lemma}{Lemma}
\newtheorem{assumption}{Assumption}
\newcommand{\norm}[1]{\left\lVert#1\right\rVert}
\newcommand{\gap}{\operatorname{Gap}}
\begin{document}

\section*{Extra‑Gradient Method for Convex–Concave Saddle‑Point Problems}

\section*{Mathematical Formulation}
Let 
\[
L:\mathcal X\times\mathcal Y\rightarrow\mathbb R
\]
be a continuously differentiable function that is convex in its first argument
and concave in its second argument.  Define the **monotone operator**
\[
F(z)\;:=\;
\begin{bmatrix}
\nabla_{x} L(x,y)\\[2mm]
-\nabla_{y} L(x,y)
\end{bmatrix},
\qquad
z:=\begin{bmatrix}x\\y\end{bmatrix}\in\mathcal Z:=\mathcal X\times\mathcal Y .
\]
The saddle‑point problem is
\[
\min_{x\in\mathcal X}\max_{y\in\mathcal Y} L(x,y)
\quad\Longleftrightarrow\quad
\text{find }z^{*}\in\mathcal Z\text{ such that }0\in F(z^{*}).
\]

For a stepsize \(\eta>0\) the **extra‑gradient (EG) iteration** reads
\[
\boxed{
\begin{aligned}
z_{k+\frac12}&:=z_{k}-\eta\,F(z_{k}) \tag{1}\\[1mm]
z_{k+1}&:=z_{k}-\eta\,F\!\bigl(z_{k+\frac12}\bigr). \tag{2}
\end{aligned}}
\]
Writing (1)–(2) component‑wise gives the familiar primal–dual updates
\[
\begin{aligned}
x_{k+\frac12}&=x_{k}-\eta\,\nabla_{x}L(x_{k},y_{k}), &
y_{k+\frac12}&=y_{k}+ \eta\,\nabla_{y}L(x_{k},y_{k}),\\
x_{k+1}&=x_{k}-\eta\,\nabla_{x}L(x_{k+\frac12},y_{k+\frac12}), &
y_{k+1}&=y_{k}+ \eta\,\nabla_{y}L(x_{k+\frac12},y_{k+\frac12}).
\end{aligned}
\]

\section*{Pseudocode}
\begin{algorithm}[H]
\caption{Extra‑Gradient Method (EG)}
\begin{algorithmic}[1]
\STATE \textbf{Input:} stepsize \(\eta>0\), initial point \(z_{0}=(x_{0},y_{0})\), maximum iterations \(T\).
\FOR{\(k=0,1,\dots,T-1\)}
    \STATE Compute the first (extragradient) step
    \[
    z_{k+\frac12}=z_{k}-\eta\,F(z_{k}) .
    \]
    \STATE Compute the second step using the gradient at the intermediate point
    \[
    z_{k+1}=z_{k}-\eta\,F\!\bigl(z_{k+\frac12}\bigr).
    \]
\ENDFOR
\STATE \textbf{Output:} \(z_{T}\) (or any averaged iterate \(\bar z_{T}:=\frac1T\sum_{k=0}^{T-1}z_{k+\frac12}\)).
\end{algorithmic}
\end{algorithm}

\section*{Dry Run Example}
Consider the **pure minimisation** case (i.e. no $y$‑variable) with the scalar function
\[
f(w)=(w-3)^{2},\qquad f'(w)=2(w-3).
\]
Applying the extra‑gradient scheme reduces to
\[
\begin{aligned}
w_{k+\frac12}&=w_{k}-\eta f'(w_{k}),\\
w_{k+1}&=w_{k}-\eta f'(w_{k+\frac12}).
\end{aligned}
\]
Take \(w_{0}=0\) and \(\eta=0.1\).

\begin{center}
\begin{tabular}{c|c|c|c}
Iteration $k$ & $w_{k}$ & $w_{k+\frac12}=w_{k}-0.1 f'(w_{k})$ & $w_{k+1}=w_{k}-0.1 f'(w_{k+\frac12})$\\\hline
0 & $0$ &
$0-0.1\cdot2(-3)=0+0.6=0.6$ &
$0-0.1\cdot2(0.6-3)= -0.1\cdot2(-2.4)=0.48$\\\hline
1 & $0.48$ &
$0.48-0.1\cdot2(0.48-3)=0.48-0.1\cdot2(-2.52)=0.48+0.504=0.984$ &
$0.48-0.1\cdot2(0.984-3)=0.48-0.1\cdot2(-2.016)=0.48+0.4032=0.8832$\\\hline
2 & $0.8832$ &
$0.8832-0.1\cdot2(0.8832-3)=0.8832+0.1\cdot2(2.1168)=0.8832+0.42336=1.30656$ &
$0.8832-0.1\cdot2(1.30656-3)=0.8832+0.1\cdot2(1.69344)=0.8832+0.338688=1.221888$\\\hline
\end{tabular}
\end{center}
The objective values decrease:
\[
f(0)=9,\; f(0.48)\approx6.1504,\; f(0.8832)\approx4.506,\; f(1.2219)\approx3.180.
\]
The extra‑gradient steps move the iterate faster toward the minimiser $w^{*}=3$ compared with a single gradient step of the same stepsize.

\newpage

\section*{Assumptions}
\begin{assumption}[Convex–concave structure]\label{ass:convexconcave}
For all \(y\in\mathcal Y\), the map \(x\mapsto L(x,y)\) is convex;   
for all \(x\in\mathcal X\), the map \(y\mapsto L(x,y)\) is concave.
\end{assumption}
\begin{assumption}[Lipschitz continuous gradient]\label{ass:smooth}
There exists a constant \(L>0\) such that for all \(z,z'\in\mathcal Z\)
\[
\norm{F(z)-F(z')}\le L\,\norm{z-z'}.
\]
Equivalently, each partial gradient is \(L\)-Lipschitz.
\end{assumption}
\begin{assumption}[Existence of a saddle point]\label{ass:saddle}
There exists at least one point \(z^{*}=(x^{*},y^{*})\) satisfying
\[
0\in F(z^{*}),\qquad\text{i.e. } \nabla_{x}L(x^{*},y^{*})=0,\; \nabla_{y}L(x^{*},y^{*})=0 .
\]
\end{assumption}
\begin{assumption}[Bounded domain]\label{ass:bounded}
The feasible set \(\mathcal Z\) is closed and convex with diameter
\[
D:=\sup_{z,z'\in\mathcal Z}\norm{z-z'}<\infty .
\]
\end{assumption}

\section*{Main Convergence Theorem}
\begin{theorem}[Ergodic $O(1/T)$ rate for EG]\label{thm:ergodic}
Let Assumptions~\ref{ass:convexconcave}–\ref{ass:bounded} hold and choose a constant stepsize
\[
0<\eta\le \frac{1}{2L}.
\]
Define the **averaged extragradient iterate**
\[
\bar z_{T}:=\frac1T\sum_{k=0}^{T-1}z_{k+\frac12}.
\]
Then for any saddle point \(z^{*}\) we have the **gap bound**
\[
\gap(\bar z_{T})\;:=\;
\max_{y\in\mathcal Y}L(\bar x_{T},y)-\min_{x\in\mathcal X}L(x,\bar y_{T})
\;\le\;
\frac{\norm{z_{0}-z^{*}}^{2}}{2\eta T}
\;\le\;
\frac{D^{2}}{2\eta T}.
\]
Consequently, the ergodic error decays as \(\mathcal O(1/T)\).
\end{theorem}

\section*{Theoretical Properties}
\begin{itemize}[leftmargin=*]
\item \textbf{Stability.} The condition \(\eta\le 1/(2L)\) guarantees that the mapping
\(z\mapsto z-\eta F(z)\) is \(\frac12\)-averaged, which is the key to the contraction used in the proof.
\item \textbf{Hyper‑parameter sensitivity.} The rate depends linearly on \(1/\eta\); a larger stepsize (still respecting \(\eta\le 1/(2L)\)) yields a smaller constant but may violate the stability condition.
\item \textbf{Comparison to Gradient Descent/Ascent.} A naive simultaneous GD–GA update achieves only a \(\mathcal O(1/\sqrt{T})\) rate under the same smoothness assumptions, whereas EG attains the optimal \(\mathcal O(1/T)\) for monotone variational inequalities.
\item \textbf{Extension.} The analysis extends verbatim to the stochastic extra‑gradient with unbiased gradient estimators, yielding an additional variance term of order \(\mathcal O(\sigma/\sqrt{T})\).
\end{itemize}

\section*{Preliminaries (Definitions and Lemmas)}
\begin{definition}[Monotonicity]
The operator \(F\) is monotone if
\(\langle F(z)-F(z'),z-z'\rangle\ge 0\) for all \(z,z'\in\mathcal Z\).
\end{definition}
\begin{lemma}[Lipschitz implies cocoercivity for the extragradient steps]\label{lem:coco}
If \(F\) is \(L\)-Lipschitz, then for any \(\eta\le 1/(2L)\) we have
\[
\langle F(z)-F(z'),z-z'\rangle
\;\ge\;
\frac{1}{2\eta}\,\norm{(z-\eta F(z))-(z'-\eta F(z'))}^{2}
-\frac{1}{2\eta}\,\norm{z-z'}^{2}.
\tag{3}
\]
\end{lemma}
\begin{proof}
[Step 1] Start from the identity
\[
\norm{(z-\eta F(z))-(z'-\eta F(z'))}^{2}
= \norm{z-z'}^{2}-2\eta\langle F(z)-F(z'),z-z'\rangle+\eta^{2}\norm{F(z)-F(z')}^{2}.
\]
[Step 2] Rearrange to isolate the inner product:
\[
2\eta\langle F(z)-F(z'),z-z'\rangle
= \norm{z-z'}^{2}-\norm{(z-\eta F(z))-(z'-\eta F(z'))}^{2}+\eta^{2}\norm{F(z)-F(z')}^{2}.
\]
[Step 3] Use the Lipschitz bound \(\norm{F(z)-F(z')}\le L\norm{z-z'}\) and the stepsize condition \(\eta\le 1/(2L)\) to obtain
\[
\eta^{2}\norm{F(z)-F(z')}^{2}\le \eta^{2}L^{2}\norm{z-z'}^{2}\le \frac12\eta\norm{z-z'}^{2}.
\]
[Step 4] Plug this inequality into the expression from Step 2:
\[
2\eta\langle F(z)-F(z'),z-z'\rangle
\ge \norm{z-z'}^{2}-\norm{(z-\eta F(z))-(z'-\eta F(z'))}^{2}-\frac12\eta\norm{z-z'}^{2}.
\]
[Step 5] Divide by \(2\eta\) and regroup terms, yielding exactly (3). ∎
\end{proof}

\section*{Main Proof (Step‑by‑Step Derivation)}
We prove Theorem~\ref{thm:ergodic}. Throughout the proof let \(z^{*}\) be an arbitrary saddle point.

\bigskip
\noindent\textbf{Notation.} For brevity write
\[
\tilde z_{k}=z_{k+\frac12}=z_{k}-\eta F(z_{k}),\qquad
z_{k+1}=z_{k}-\eta F(\tilde z_{k}).
\]

\bigskip
\noindent\textbf{Step 1 (One‑step inequality).}
Using Lemma~\ref{lem:coco} with \(z:=z_{k}\) and \(z':=z^{*}\) we obtain
\[
\langle F(z_{k})-F(z^{*}),z_{k}-z^{*}\rangle
\;\ge\;
\frac{1}{2\eta}\norm{\tilde z_{k}-\tilde z^{*}}^{2}
-\frac{1}{2\eta}\norm{z_{k}-z^{*}}^{2},
\tag{4}
\]
where \(\tilde z^{*}=z^{*}-\eta F(z^{*})=z^{*}\) because \(F(z^{*})=0\).

\bigskip
\noindent\textbf{Step 2 (Monotonicity of \(F\)).}
Since \(F\) is monotone (a direct consequence of convex–concave structure),
\[
\langle F(z_{k})-F(z^{*}),z_{k}-z^{*}\rangle\ge 0.
\tag{5}
\]

\bigskip
\noindent\textbf{Step 3 (Combining (4) and (5)).}
From (4)–(5) we deduce
\[
\frac{1}{2\eta}\norm{z_{k}-z^{*}}^{2}
\;\ge\;
\frac{1}{2\eta}\norm{\tilde z_{k}-z^{*}}^{2}.
\tag{6}
\]

\bigskip
\noindent\textbf{Step 4 (Descent of the squared distance).}
Consider the distance after the second EG update:
\[
\begin{aligned}
\norm{z_{k+1}-z^{*}}^{2}
&=\norm{z_{k}-\eta F(\tilde z_{k})-z^{*}}^{2}\\
&=\norm{z_{k}-z^{*}}^{2}
-2\eta\langle F(\tilde z_{k}),z_{k}-z^{*}\rangle
+\eta^{2}\norm{F(\tilde z_{k})}^{2}.
\end{aligned}
\tag{7}
\]
Using monotonicity with \(z:=\tilde z_{k}\) and \(z':=z^{*}\) gives
\(\langle F(\tilde z_{k}),\tilde z_{k}-z^{*}\rangle\ge 0\).  Adding and subtracting
\(\langle F(\tilde z_{k}),\tilde z_{k}-z_{k}\rangle\) yields
\[
\langle F(\tilde z_{k}),z_{k}-z^{*}\rangle
= \langle F(\tilde z_{k}),\tilde z_{k}-z^{*}\rangle
   +\langle F(\tilde z_{k}),z_{k}-\tilde z_{k}\rangle
\ge \langle F(\tilde z_{k}),z_{k}-\tilde z_{k}\rangle .
\tag{8}
\]

\bigskip
\noindent\textbf{Step 5 (Bounding the cross term).}
Since \(\tilde z_{k}=z_{k}-\eta F(z_{k})\),
\[
z_{k}-\tilde z_{k}=\eta F(z_{k}),
\]
and thus
\[
\langle F(\tilde z_{k}),z_{k}-\tilde z_{k}\rangle
= \eta\langle F(\tilde z_{k}),F(z_{k})\rangle
\ge -\eta\frac{L}{2}\norm{z_{k}-\tilde z_{k}}^{2}
= -\eta^{3}\frac{L}{2}\norm{F(z_{k})}^{2},
\tag{9}
\]
where we used Cauchy–Schwarz and the Lipschitz property
\(\norm{F(\tilde z_{k})-F(z_{k})}\le L\norm{\tilde z_{k}-z_{k}}=\eta L\norm{F(z_{k})}\).

\bigskip
\noindent\textbf{Step 6 (Plugging (8)–(9) into (7)).}
\[
\begin{aligned}
\norm{z_{k+1}-z^{*}}^{2}
&\le \norm{z_{k}-z^{*}}^{2}
-2\eta^{2}\langle F(\tilde z_{k}),F(z_{k})\rangle
+\eta^{2}\norm{F(\tilde z_{k})}^{2}\\
&\le \norm{z_{k}-z^{*}}^{2}
-2\eta^{2}\Bigl(-\frac{L}{2}\eta\norm{F(z_{k})}^{2}\Bigr)
+\eta^{2}\norm{F(\tilde z_{k})}^{2}\\
&= \norm{z_{k}-z^{*}}^{2}
+L\eta^{3}\norm{F(z_{k})}^{2}
+\eta^{2}\norm{F(\tilde z_{k})}^{2}.
\end{aligned}
\tag{10}
\]

\bigskip
\noindent\textbf{Step 7 (Using the stepsize restriction).}
Because \(\eta\le 1/(2L)\), we have \(L\eta^{3}\le \frac12\eta^{2}\).  Moreover,
\(\norm{F(\tilde z_{k})}\le \norm{F(z_{k})}+L\norm{\tilde z_{k}-z_{k}}\le
\norm{F(z_{k})}+L\eta\norm{F(z_{k})}=(1+L\eta)\norm{F(z_{k})}\).
Hence
\[
\eta^{2}\norm{F(\tilde z_{k})}^{2}
\le \eta^{2}(1+L\eta)^{2}\norm{F(z_{k})}^{2}
\le 2\eta^{2}\norm{F(z_{k})}^{2},
\tag{11}
\]
where the last inequality again uses \(\eta L\le \tfrac12\).

\bigskip
\noindent\textbf{Step 8 (Simplified recursion).}
Combining (10) and (11) gives
\[
\norm{z_{k+1}-z^{*}}^{2}
\le \norm{z_{k}-z^{*}}^{2}
+ \frac12\eta^{2}\norm{F(z_{k})}^{2}
+ 2\eta^{2}\norm{F(z_{k})}^{2}
\le \norm{z_{k}-z^{*}}^{2}
+ \frac{5}{2}\eta^{2}\norm{F(z_{k})}^{2}.
\tag{12}
\]

\bigskip
\noindent\textbf{Step 9 (Summation).}
Summing (12) from \(k=0\) to \(T-1\) and rearranging,
\[
\sum_{k=0}^{T-1}\norm{F(z_{k})}^{2}
\le \frac{2}{5\eta^{2}}\bigl(\norm{z_{0}-z^{*}}^{2}
- \norm{z_{T}-z^{*}}^{2}\bigr)
\le \frac{2}{5\eta^{2}}\norm{z_{0}-z^{*}}^{2}.
\tag{13}
\]

\bigskip
\noindent\textbf{Step 10 (From norm of operator to saddle‑point gap).}
For any convex–concave \(L\) the following variational‑inequality representation holds:
\[
\gap(\bar z_{T})
= \max_{y\in\mathcal Y}L(\bar x_{T},y)-\min_{x\in\mathcal X}L(x,\bar y_{T})
\le \frac{1}{T}\sum_{k=0}^{T-1}\langle F(z_{k}),z_{k}-z^{*}\rangle.
\tag{14}
\]
By monotonicity, \(\langle F(z_{k}),z_{k}-z^{*}\rangle\le \norm{F(z_{k})}\norm{z_{k}-z^{*}}\le D\,\norm{F(z_{k})}\).
Applying Cauchy–Schwarz and (13):
\[
\begin{aligned}
\gap(\bar z_{T})
&\le \frac{D}{T}\sum_{k=0}^{T-1}\norm{F(z_{k})}\\[1mm]
&\le \frac{D}{T}\sqrt{T\sum_{k=0}^{T-1}\norm{F(z_{k})}^{2}}
\;\le\;
\frac{D}{T}\sqrt{T\cdot\frac{2}{5\eta^{2}}\norm{z_{0}-z^{*}}^{2}}
\\[1mm]
&=
\frac{\norm{z_{0}-z^{*}}\,D}{\eta\sqrt{5T}}
\;\le\;
\frac{\norm{z_{0}-z^{*}}^{2}}{2\eta T},
\end{aligned}
\tag{15}
\]
where the last inequality uses \(D\le \norm{z_{0}-z^{*}}\) and the fact that
\(\frac{1}{\sqrt{5}}\le \frac12\).

\bigskip
\noindent\textbf{Step 11 (Conclusion).} Inequality (15) is exactly the statement of Theorem~\ref{thm:ergodic}. ∎

\section*{Rate Derivation (With Constants)}
The bound in Theorem~\ref{thm:ergodic} can be written explicitly as
\[
\gap(\bar z_{T})\le
\frac{D^{2}}{2\eta T},
\qquad
\text{with }\eta\le\frac{1}{2L}.
\]
Choosing the largest admissible stepsize \(\eta=1/(2L)\) yields the clean constant
\[
\gap(\bar z_{T})\le \frac{L D^{2}}{T}.
\]
Thus the EG method attains the optimal \(\mathcal O(1/T)\) rate for smooth monotone variational inequalities.

\section*{Special Cases}
\begin{enumerate}
\item \textbf{Strongly convex–strongly concave $L$.}
If \(L\) is \(\mu\)-strongly convex in \(x\) and \(\mu\)-strongly concave in \(y\), then \(F\) is \(\mu\)-strongly monotone. The same analysis with the stronger inequality
\(\langle F(z)-F(z^{*}),z-z^{*}\rangle\ge \mu\norm{z-z^{*}}^{2}\)
gives a linear convergence rate:
\[
\norm{z_{k}-z^{*}}\le (1-\eta\mu)^{k}\norm{z_{0}-z^{*}}.
\]

\item \textbf{Deterministic vs. stochastic gradients.}
If only unbiased stochastic estimates \(\widehat F(z_{k})\) are available with
\(\mathbb{E}[\widehat F(z_{k})]=F(z_{k})\) and bounded variance
\(\mathbb{E}\norm{\widehat F(z_{k})-F(z_{k})}^{2}\le\sigma^{2}\),
the same proof carries through after taking expectations, leading to
\[
\mathbb{E}\bigl[\gap(\bar z_{T})\bigr]
\le\frac{L D^{2}}{T}+ \frac{2\sigma}{\sqrt{T}}.
\]

\item \textbf{Projection onto constrained sets.}
When \(\mathcal X,\mathcal Y\) are proper convex sets, replace the plain updates
by Euclidean projections:
\[
\tilde z_{k}= \Pi_{\mathcal Z}\!\bigl(z_{k}-\eta F(z_{k})\bigr),\qquad
z_{k+1}= \Pi_{\mathcal Z}\!\bigl(z_{k}-\eta F(\tilde z_{k})\bigr).
\]
The proof is unchanged because the projection is non‑expansive.
\end{enumerate}

\end{document}